\documentclass{beamer}

\usepackage[spanish]{babel}
\usepackage[utf8]{inputenc}
\usepackage[T1]{fontenc}
\usepackage{listings}


\usetheme{Madrid} % Puedes cambiar el tema según prefieras

\title{Gestión de Usuarios, Grupos y Permisos en Ubuntu}
\author{Tu Nombre}
\date{\today}

\begin{document}

%-------------------- PORTADA --------------------
\begin{frame}
    \titlepage
\end{frame}

%-------------------- ÍNDICE ---------------------
\begin{frame}{Contenido}
    \tableofcontents
\end{frame}

%=================================================
\section{Ejemplos Reales}
%=================================================

%-------------------- Slide 13 --------------------
\begin{frame}[fragile]{Administración de un Servidor Web}
    \begin{enumerate}
        \item Crear grupo \texttt{webdev}:
        \begin{lstlisting}
sudo groupadd webdev
        \end{lstlisting}
        \item Agregar usuarios al grupo:
        \begin{lstlisting}
sudo usermod -aG webdev juan
        \end{lstlisting}
        \item Asignar grupo y permisos en /var/www/html/:
        \begin{lstlisting}
sudo chown -R root:webdev /var/www/html/
sudo chmod -R 775 /var/www/html/
        \end{lstlisting}
    \end{enumerate}
\end{frame}

%-------------------- Slide 14 --------------------
\begin{frame}{Directorio Compartido}
    \begin{enumerate}
        \item Crear grupo \texttt{compartidos}:
        \begin{lstlisting}
sudo groupadd compartidos
        \end{lstlisting}
        \item Agregar usuarios al grupo:
        \begin{lstlisting}
sudo usermod -aG compartidos ana
        \end{lstlisting}
        \item Crear directorio y asignar permisos:
        \begin{lstlisting}
sudo mkdir /home/compartido
sudo chown -R root:compartidos /home/compartido
sudo chmod -R 770 /home/compartido
        \end{lstlisting}
    \end{enumerate}
\end{frame}

%=================================================
\section{Ejercicios Prácticos}
%=================================================

%-------------------- Slide 15 --------------------
\begin{frame}{Ejercicio 1: Creación de un Usuario}
    \textbf{Objetivo:} Crear un usuario \texttt{prueba1} y verificarlo.
    \begin{enumerate}
        \item \texttt{sudo adduser prueba1}
        \item Comprobar: \texttt{id prueba1}
    \end{enumerate}
    \textbf{Resultado esperado:} \texttt{uid=1001(prueba1) gid=1001(prueba1) ...}
\end{frame}

%-------------------- Slide 16 --------------------
\begin{frame}{Ejercicio 2: Creación de un Grupo y Asignación de Usuario}
    \begin{enumerate}
        \item \texttt{sudo groupadd practica}
        \item \texttt{sudo usermod -aG practica prueba1}
    \end{enumerate}
    \textbf{Comprobación:} \texttt{groups prueba1} \\
    \textbf{Esperado:} \texttt{prueba1 : prueba1 practica}
\end{frame}

%-------------------- Slide 17 --------------------
\begin{frame}{Ejercicio 3: Cambio de Propietario de un Archivo}
    \begin{enumerate}
        \item Crear \texttt{ejemplo.txt} en tu \textbf{home}.
        \item Cambiar propietario a \texttt{prueba1}.
    \end{enumerate}
    \textbf{Comprobación:} \texttt{ls -l ejemplo.txt} \\
    \textbf{Esperado:} \texttt{-rw-r--r-- 1 prueba1 usuario ... ejemplo.txt}
\end{frame}

%-------------------- Slide 18 --------------------
\begin{frame}{Ejercicio 4: Cambiar Permisos de un Archivo}
    \begin{enumerate}
        \item Crear \texttt{script.sh}.
        \item Dar permisos de ejecución a todos (\texttt{chmod 755 script.sh}).
    \end{enumerate}
    \textbf{Comprobación:} \texttt{ls -l script.sh} \\
    \textbf{Esperado:} \texttt{-rwxr-xr-x ... script.sh}
\end{frame}

%-------------------- Slide 19 --------------------
\begin{frame}{Ejercicio 5: Directorio Compartido con Grupo}
    \begin{enumerate}
        \item \texttt{sudo groupadd equipo1}
        \item Crear \texttt{usuario2} y añadirlo a \texttt{equipo1}.
        \item Directorio \texttt{/home/compartido/} con grupo \texttt{equipo1}.
        \item \texttt{chmod -R 770 /home/compartido/}
    \end{enumerate}
    \textbf{Comprobación:} \texttt{ls -ld /home/compartido/} \\
    \textbf{Esperado:} \texttt{drwxrwx--- ... root equipo1 ...}
\end{frame}

%-------------------- Slide 20 --------------------
\begin{frame}{Ejercicio 6: Configurar el Sticky Bit}
    \begin{enumerate}
        \item Crear \texttt{/home/publico/} con permisos de escritura para todos.
        \item Activar Sticky Bit: \texttt{chmod +t /home/publico/}
    \end{enumerate}
    \textbf{Comprobación:} \texttt{ls -ld /home/publico/} \\
    \textbf{Esperado:} \texttt{drwxrwxrwt 2 root root ... /home/publico/}
\end{frame}

%-------------------- Slide 21 --------------------
\begin{frame}{Ejercicio 7: Usuario con Shell Restringido}
    \begin{enumerate}
        \item Crear \texttt{limitado} con shell \texttt{/bin/rbash}.
        \item Probar \texttt{cd /} y ver error.
    \end{enumerate}
    \textbf{Comprobación:} \texttt{rbash: cd: restricted}
\end{frame}

%-------------------- Slide 22 --------------------
\begin{frame}{Ejercicio 8: Asignar Permisos con SetGID}
    \begin{enumerate}
        \item \texttt{sudo groupadd proyectos}
        \item Crear \texttt{/home/proyectos/} y asignar grupo.
        \item \texttt{chmod g+s /home/proyectos/}
    \end{enumerate}
    \textbf{Comprobación:} \texttt{ls -ld /home/proyectos/} \\
    \textbf{Esperado:} \texttt{drwxrws--- ... root proyectos ...}
\end{frame}

%=================================================
\section{Conclusiones}
%=================================================

%-------------------- Slide 23 --------------------
\begin{frame}{Conclusión}
    \begin{block}{Resumen}
    \begin{itemize}
        \item Administrar usuarios, grupos y permisos es esencial para la seguridad y eficiencia en Ubuntu.
        \item Comprender los \textbf{permisos (r, w, x)}, la asignación de propietarios y grupos permite un control preciso.
        \item Los ejercicios propuestos te ayudan a practicar y verificar tu aprendizaje de forma directa.
    \end{itemize}
    \end{block}
    \vspace{0.5cm}
    \centering
    \alert{¡Si tus resultados no coinciden con los esperados, revisa y repite los pasos!}
\end{frame}

%-------------------- FIN --------------------
\end{document}
